\documentclass[11pt,a4paper]{article}

\usepackage[utf8]{inputenc}
\usepackage[T1]{fontenc}
\usepackage[french]{babel}
\usepackage[margin=2.5cm]{geometry}
\usepackage{listings}
\usepackage{xcolor}
\usepackage{enumitem}
\usepackage{fancyhdr}
\usepackage{titlesec}

\definecolor{bleu}{RGB}{21,67,140}

\lstdefinestyle{min}{
    basicstyle=\ttfamily\footnotesize,
    breaklines=true,
    frame=none,
    numbers=none,
    tabsize=4,
    language=C++,
    backgroundcolor=\color{white}
}
\lstset{style=min}

\setlist[itemize]{leftmargin=1.3em, itemsep=2pt, topsep=3pt}
\setlist[enumerate]{leftmargin=1.5em, itemsep=2pt, topsep=3pt}

\pagestyle{fancy}
\fancyhf{}
\renewcommand{\headrulewidth}{0.4pt}
\fancyhead[L]{\small\color{bleu}ERRAMI Youssef}
\fancyhead[R]{\small\color{bleu}TP1 -- StackDLL}
\fancyfoot[C]{\thepage}

\titleformat{\section}{\normalfont\large\bfseries\color{bleu}}{\thesection}{0.8em}{}
\titlespacing*{\section}{0pt}{12pt}{6pt}
\titleformat{\subsection}{\normalfont\bfseries\color{bleu}}{\thesubsection}{0.8em}{}
\titlespacing*{\subsection}{0pt}{8pt}{4pt}

\newcommand{\pt}[1]{\textbf{\color{bleu}#1}}

\begin{document}

{\Large\bfseries\color{bleu} Compilation et DLL -- ce qu'on a fait}\\[2pt]
\small ERRAMI Youssef -- ILISI 2, C++ Avancé, TP1

\vspace{4pt}\hrule\vspace{10pt}

\section{Les 4 étapes de compilation}

\begin{itemize}
\item \pt{Prétraitement} : remplace les \texttt{\#include}, \texttt{\#define}, \texttt{\#ifdef} par leur contenu. $\to$ un seul fichier texte.
\item \pt{Compilation} : traduit le C++ en \pt{assembleur}. Les classes/templates/boucles disparaissent au profit de \texttt{mov}, \texttt{cmp}, \texttt{jmp}, \texttt{call}.
\item \pt{Assemblage} : traduit l'assembleur en binaire $\to$ fichier \texttt{.obj}. Pas encore exécutable (symboles non résolus). \textit{C'est ici que MinGW a bloqué (section~\ref{sec:dg}).}
\item \pt{Édition de liens} : \texttt{link.exe}/\texttt{ld} relie les \texttt{.obj} entre eux $\to$ \texttt{.exe} ou \texttt{.dll}.
\end{itemize}

\section{DLL vs bibliothèque statique}

\begin{itemize}
\item \pt{Statique} (\texttt{.lib}/\texttt{.a}) : le code est \textbf{copié} dans l'exécutable à la compilation.
\item \pt{Dynamique} (\texttt{.dll}) : le code reste un fichier séparé, chargé \textbf{au lancement}.
\end{itemize}

Pour que le même \texttt{Stack.h} serve à fabriquer la DLL \textbf{et} à l'utiliser, une seule macro :

\begin{lstlisting}[language=C++]
#ifdef STACKDLL_EXPORTS
  #define STACKDLL_API __declspec(dllexport)   // on fabrique la DLL
#else
  #define STACKDLL_API __declspec(dllimport)   // on l'utilise
#endif
\end{lstlisting}

\section{Les fichiers générés}

Commande : \texttt{cl /LD Stack.cpp Expression.cpp /Fe:StackDLL.dll} $\to$ 3 fichiers, pas 1.

\begin{itemize}
\item \pt{.obj} (un par .cpp) : code machine + symboles réclamés/offerts. Pas exécutable seul.
\item \pt{.dll} : le binaire final + une table d'export (liste des symboles accessibles).
\item \pt{.lib} (import library) : \textbf{aucun code machine} dedans ! Juste une table \og ce symbole, va le chercher dans la DLL \fg. Sert uniquement à lier \texttt{main.cpp}.
\item \pt{.exp} : intermédiaire temporaire du link, sert à fabriquer le \texttt{.dll} et le \texttt{.lib} ensemble. Jetable après.
\item \pt{.exe} : contient \texttt{main.cpp}, pas le code des piles. Sans la DLL à côté, il ne démarre pas.
\end{itemize}

\texttt{dumpbin /exports StackDLL.dll} $\to$ 38 symboles, du type :
\begin{lstlisting}
?push@?$DynamicStack@D@@QEAAXAEBD@Z
\end{lstlisting}
Nom illisible = \pt{name mangling} : le compilateur encode classe + type du template + paramètres dans le nom, pour distinguer \texttt{push<char>} de \texttt{push<int>}.

\section{Lire l'assembleur généré}

Commande : \texttt{cl /c /FAs Stack.cpp} $\to$ vrai fichier \texttt{Stack.asm}. Extrait pour \texttt{isEmpty()} (= \texttt{return topIndex == -1;}) :

\begin{lstlisting}
mov  rax, QWORD PTR this$[rsp]
cmp  DWORD PTR [rax+12], -1    ; topIndex == -1 ?
jne  SHORT $LN3
mov  DWORD PTR tv66[rsp], 1
jmp  SHORT $LN4
$LN3: mov DWORD PTR tv66[rsp], 0
$LN4: movzx eax, BYTE PTR tv66[rsp]
\end{lstlisting}

On reconnaît le \texttt{if}/\texttt{else} du code, écrit avec des sauts au lieu de mots-clés.

\pt{Peut-on retraduire ça en \texttt{Stack.cpp} ?} Partiellement :
\begin{itemize}
\item les noms de variables disparaissent (\texttt{topIndex} $\to$ juste \og offset 12 \fg) ;
\item les classes/templates sont aplatis en blocs de code séparés ;
\item avec optimisation (\texttt{/O2}), l'ordre des instructions change et ne colle plus au source.
\end{itemize}
Des outils existent pour aller plus loin : désassembleurs (\texttt{objdump}, \texttt{dumpbin /disasm}) et décompilateurs (Ghidra, IDA) -- jamais une reconstruction exacte, juste une approximation à relire.

\section{Le blocage Device Guard}
\label{sec:dg}

Ce qu'on a vraiment rencontré pendant le TP :

\begin{itemize}
\item MinGW (\texttt{g++}) échouait sans message. En creusant (\texttt{g++ -v}) : ça plantait à l'assemblage.
\item Message réel en appelant \texttt{as.exe} directement : \textit{``An Application Control policy has blocked this file''}.
\item Avec MSVC, la compilation réussit... mais \textbf{lancer} \texttt{main.exe} donne : \textit{``blocked by your organization's Device Guard policy''}.
\end{itemize}

\pt{Ce que ça veut dire} : Device Guard (WDAC) bloque au niveau machine tout exécutable non approuvé, peu importe sa validité. Vérifié avec \texttt{dumpbin} : les 38 symboles sont bien là, le binaire est correct -- c'est juste le \textbf{droit de lancer} qui est refusé.

\pt{À retenir} : ``ça compile et se lie'' et ``ça peut s'exécuter sur cette machine'' sont deux questions indépendantes.

\section{Questions probables}

\begin{enumerate}[label=\pt{Q\arabic*.}]
\item \textbf{Pourquoi \texttt{extern template class DynamicStack<int>;} ?} \\
Évite que chaque \texttt{.cpp} réinstancie sa propre copie du template. Une seule version, celle de la DLL.

\item \textbf{Pourquoi encadrer ça avec \texttt{\#ifndef STACKDLL\_EXPORTS} ?} \\
Sans ça, MSVC renvoie l'avertissement C4910 : on ne peut pas dire \og j'exporte \fg{} et \og c'est externe \fg{} en même temps dans le même fichier.

\item \textbf{Le \texttt{.lib} d'une DLL = un \texttt{.lib} statique classique ?} \\
Non. Le statique contient du code. L'import library ne contient que des symboles.

\item \textbf{Si on efface \texttt{StackDLL.dll}, que se passe-t-il ?} \\
\texttt{main.exe} ne démarre pas (\og DLL introuvable \fg). La résolution se fait au lancement, pas à la compilation.

\item \textbf{Changer de compilateur évite le blocage Device Guard ?} \\
Non pour le lancement : n'importe quel \texttt{.exe} non signé est bloqué, MSVC ou MinGW pareil. (Mais changer de compilateur a réglé le blocage à l'assemblage.)

\item \textbf{Pourquoi les noms exportés sont illisibles ?} \\
Name mangling : le C++ autorise la surcharge, donc le nom binaire doit encoder le type pour que l'éditeur de liens distingue les versions.
\end{enumerate}

\section{En résumé}

\begin{enumerate}
\item Un \texttt{.obj} = un fragment de code, pas un programme.
\item Une DLL = 3 fichiers (\texttt{.dll}, \texttt{.lib} d'import, \texttt{.exp} jetable). Le \texttt{.lib} ne contient pas de code.
\item Compiler/lier avec succès $\neq$ avoir le droit de s'exécuter sur la machine.
\end{enumerate}

\end{document}
